\documentclass{beamer}
%\usepackage[english]{babel}
%\usepackage{graphicx}
\usepackage{bm}

%%%%%%%%%%%%%%%%%%%%%%%%%%%%%%%%%%%%%%%%%%%%%%%%%%%%%%%%%%%%%%%%%%%%%%%%%%%%%%%%

\title{\Large PrintConcrete: Isogeometric FEM simulator \newline
       for concrete 3D printing}
\author{Knut Morten Okstad$^{\dag}$}

\institute{$^\dag$ SINTEF Digital, Trondheim, Norway}

\hypersetup{colorlinks=true,urlcolor=blue}
\hypersetup{pdfauthor=Knut Morten Okstad}
\hypersetup{pdfsubject=Overal description of the PrintConrete IFEM application}
\hypersetup{pdfkeywords=FEM IGA Nonlinear-elasticity 3D-printing Concrete}

\begin{document} %%%%%%%%%%%%%%%%%%%%%%%%%%%%%%%%%%%%%%%%%%%%%%%%%%%%%%%%%%%%%%%

\frame{\titlepage}

\frame{\frametitle{The PrintConcrete IFEM simulator}
  \begin{itemize}
  \item Quasi-static nonlinear FE simulator of the elasticity PDE
    $$
    \nabla\bm{\sigma} + \bm{f} = \bm{0}
    \quad\mbox{where}\quad \bm{\sigma} = \bm{C}(t):\bm{\varepsilon}(\bm{u})
    $$
    where $\bm{C}$ represents a time-dependent material tensor, \newline
    and $\bm{f}$ consists of the gravity load.
  \item The printing process is simulated by an element-birth feature in the solver,
    i.e., the stiffness and and gravity loads for an element is set to zero
    until the birth time of the element.
  \item Material properties {\small (Young's modulus $E$ and yield stress limit $\sigma_y$)}
    depend on the element age (time since element birth).
  \item The nonlinear PDE is linarized based on an Updated Lagrange formulation,
    and solved using Newton-Raphson iterations.% on each load increment.
  \item Isogeometric analysis: Quadratic or cubic splines are used as basis functions
    in the FE implementation.
  \item The simulator is built on the IFEM-Elasticiy package \url{https://github.com/OPM/IFEM-Elasticity}.
  \end{itemize}}

\end{document} %%%%%%%%%%%%%%%%%%%%%%%%%%%%%%%%%%%%%%%%%%%%%%%%%%%%%%%%%%%%%%%%%
